\documentclass[11pt,a4paper]{article}

\usepackage[margin=1in]{geometry}
\usepackage{parskip}
\usepackage{titlesec}
\usepackage{enumitem}
\usepackage{hyperref}

\titleformat{\section}{\large\bfseries}{\thesection.}{0.5em}{}
\titleformat{\subsection}{\normalsize\bfseries}{\thesubsection}{0.5em}{}
\setlist{itemsep=2pt, topsep=4pt}

\hypersetup{colorlinks=true, linkcolor=black, urlcolor=blue}

\newcommand{\doctitle}[2]{
  \begin{center}
    {\LARGE\bfseries #1}\\[4pt]
    {\large #2}\\[6pt]
    Sahithyan Kandathasan \\
    Index No: 230557T
  \end{center}
  \vspace{1em}
}


\begin{document}

\doctitle{Business Proposal}{Jaffna SeaWeave: Cultivated Sea Moss for Food and Cosmetic Markets}

\section{Introduction and Resource}
Jaffna SeaWeave is a proposed small business that cultivates and processes sea moss (red algae, primarily \textit{Eucheuma} and \textit{Gracilaria} species) in the shallow lagoon waters around the Jaffna peninsula and its islands, such as Delft and Analaitivu. These lagoons offer the calm currents, warm temperatures, and salinity needed for sea moss to grow on suspended cultivation lines, and the crop regrows from cuttings, so a single initial stock can sustain repeated harvests. Demand for sea moss as a health food, thickener, and cosmetic ingredient has grown internationally, met mostly through imports, while Jaffna's lagoons are largely unused for this purpose.

\section{Business Idea and Production}
The business will farm sea moss using the rope/raft method: seedling cuttings tied to floating lines, grown for six to eight weeks, then harvested by hand, washed, sun-dried, and sorted. Output will initially be sold as raw dried sea moss; as the business grows, it can add simple processing into sea moss gel to capture more value. The production process is:

\begin{center}
Seedling stock $\rightarrow$ Line/raft cultivation $\rightarrow$ Harvesting $\rightarrow$ Washing $\rightarrow$ Sun-drying $\rightarrow$ Sorting $\rightarrow$ Packaging $\rightarrow$ Distribution
\end{center}

Drying quality and sorting matter directly to price, since moisture content and cleanliness affect grade in both food and cosmetic markets.

\section{Objectives}
\begin{itemize}
  \item Establish a locally grown, sustainably farmed source of sea moss from Jaffna's lagoon waters.
  \item Supply domestic and export demand for sea moss as a food and cosmetic ingredient.
  \item Create income for coastal households as farm labour or seedling suppliers.
  \item Build a regional brand around Jaffna-grown sea moss.
\end{itemize}

\section{Products and Target Market}
The initial product is dried sea moss, sold in bulk and retail packaging; a sea moss gel can follow once processing capacity and food-safety certification are in place. Customers include local health-food and cosmetic manufacturers, wholesale buyers supplying Colombo, and, as the business scales, export buyers in markets with established sea moss demand.

\section{Monopoly Market Structure and Barriers to Entry}
For the economics activity, the business is considered within a narrowly defined market: cultivated sea moss sourced specifically from Jaffna's lagoon waters, not all seaweed or algae products generally. As the first organized cultivator in the area, the business could secure long-term aquaculture leases over the best sites and build cultivation know-how and buyer relationships before competitors enter, creating a dominant position within this niche.

The resulting barriers to entry include the limited number of suitable lagoon sites, the cost and permitting needed to secure a lease, the time needed to build cultivation know-how (poor handling causes crop loss), and an early mover's existing buyer relationships and brand recognition. These barriers do not block competition in the broader seaweed market, but they make entry into this specific niche harder for later entrants.

\section{Pricing, Costs, and Marketing}
Revenue comes mainly from dried sea moss and, later, gel sales. Because buyers can source sea moss from other countries, pricing has to account for import prices and quality requirements alongside production costs and margin. Major costs include seedling stock, cultivation lines and rafts, lease fees, labour, drying and storage infrastructure, packaging, food-grade certification, and transport; lagoon space itself is cheap relative to land-based farming, but labour and drying infrastructure are ongoing costs.

Marketing will emphasise a traceable local source and quality suited to both food and cosmetic use, reaching buyers through health-food retailers, direct outreach to cosmetic manufacturers, and local trade fairs.

\section{Environmental Benefits and Risks}
Sea moss needs no freshwater, fertiliser, or land, and can improve water quality by absorbing excess nutrients, while giving coastal households income without competing with fishing grounds. Key risks are storm damage to cultivation lines, disease or biofouling, inconsistent drying due to weather, and delays in food-grade certification for export. These can be reduced through site diversification, staff training, investment in covered drying as the business grows, and early engagement with certification requirements.

\section{Conclusion}
Jaffna SeaWeave proposes a low-input aquaculture business using a natural resource unique to the Jaffna peninsula's lagoon waters to supply a growing ingredient market. Its monopoly characteristics apply within a clearly defined local niche, supported by early access to lagoon sites, cultivation know-how, and buyer relationships built ahead of later entrants.

\end{document}
